\begin{questionS}
We require a concurrent object for asynchronously distributing messages
between |n| external clients, each of which has an identity in the range
$\interval{0}{\sm n}$.  The object should have the following signature.
\begin{scala}
class Distribution(n: Int){
  def send(s: Int, r: Int, m: String): Unit = ...
  def receive(r: Int): (Int, String) = ...
  def shutdown(): Unit = ...
}
\end{scala}
%
Client~|s| can send message~|m| to client~|r| (where maybe $\sm r = \sm s$) by
calling |send(s, r, m)|.  If subsequently client~|r| calls |receive(r)| then
it can receive the pair \SCALA{(s, m)} (or a pair corresponding to another
call of |send|).  The |shutdown| operation shuts down the system, as usual. 

Implement this class.  Internally, you should use a ring of~|n| nodes to pass
messages.  Each node needs to be able to handle tokens passed on the ring, or
relevant calls of |send| or |receive|.  You will need to think about how to
avoid deadlocks.

Test your code.
\end{questionS}

%%%%%%%%%%%%%%%%%%%%%%%%%%%%%%%%%%%%%%%%%%%%%%%%%%%%%%%

\begin{answerS}
My code is below, explained by comments.  Tokens are passed round the ring;
each token includes the sender, receiver and contents of a message.  Each node
uses an alternation to handle both tokens and messages corresponding to new
|send|s; the helper function |dispatch| either passes a token, or sends to a
corresponding |receive| operation.
%
\begin{scala}
class Distribution(n: Int){
  /** A token is a triple £(s,r,m)£, indicating message £m£ being sent from £s£ to £r£. */
  type Token = (Int, Int, String)
  /** Channels for the ring.  £ring(i)£ goes from node £(i-1) mod n£ to node £i£;
    * so node £i£ inputs on £ring(i)£ and outputs on £chan((i+1)\%n)£. */  
  private val ring = Array.fill(n)(new UnboundedBuffChan[Token])

  /** Input channels into the ring. */
  private val in = Array.fill(n)(new UnboundedBuffChan[(Int, String)])

  /** Output channels from the ring. */
  private val out = Array.fill(n)(new UnboundedBuffChan[(Int, String)])

  /** Client £s£ sends message £m£ for thread £r£. */
  def send(s: Int, r: Int, m: String) = {
    require(0 <= s && s < n && 0 <= r && r < n); in(s)!((r,m))
  }

  /** Client £r£ receives a message, together with the identity of the sending
    * client. */
  def receive(r: Int): (Int, String) = { require(0 <= r && r < n); out(r)?() }

  /** Ring node with identity £me£. */
  private def node(me: Int) = thread(s"node $me"){
    val left = ring(me); val right = ring((me+1)%n)
    /* Deal with token (s,r,m), */
    def dispatch(s: Int, r: Int, m: String) = {
      if(r == me) out(me)!(s,m) else right!(s,r,m)
    }
    serve(
      left =?=> { case (s,r,m) => dispatch(s,r,m) }
      | in(me) =?=> { case (r,m) => dispatch(me,r,m) }
    )
  }

  // Fork off the node threads. */
  fork(|| (for(i <- 0 until n) yield node(i)))

  /** Shut down the system. */
  def shutdown() = {
    ring.foreach(_.close()); in.foreach(_.close()); out.foreach(_.close())
  }
}
\end{scala}
% 
It is necessary to use unbounded buffering on the |ring| channels.  In
principle, more and more tokens could be added to the system if clients call
|send| sufficiently often (although this seems unlikely in practice).  We also
use unbounded buffering for the |out| channels, to avoid blocking the ring if
a client does not call |receive|, and likewise on the |in| channels, to avoid
blocking |send| operations.

For testing, we arrange for each client to send |k| messages to each client,
including itself, in a random order.  Thus each client sends and receives
$\sm{iters} = \sm{n} \times \sm{k}$ messages.  Each client sends |String|s
representing $\interval{0}{\sm{iters}}$, in order.  Each client stores
information about what it receives.  For convenience, we implement each client
as a pair of threads, to handle |send|s and |receive|s, respectively.

When the clients terminate we check that if client~|r| received message~|m|
from~|s|, then |s| sent~|m| to~|r|.  Assuming the system terminates, this
implies the converse. 

The following function performs a single test. 
%
\begin{scala}
  def doTest() = {
    // We use £n£ workers.  Each sends £k£ values to each worker (including
    // itself).  So each thread performs a total of £iters£ sends and receives.
    // Each thread sends values £[0..iters)£ in order, converted to £String£s.
    val n = 5; val k = 4; val iters = n*k
    val dist = new Distribution(n)
    // If £sends(s)(i) = r£ then £sender(s)£'s £i£'th send is to client £r£ with
    // value £i.toString£.
    val sends = new Array[Array[Int]](n)
    // If £receives(r)(i) = (s,m)£, then £receiver(r)£'s £i£'th receive was £(s,m)£.
    val receives = Array.ofDim[(Int,String)](n,iters)
    def sender(me: Int) = thread(s"sender $me"){
      // This worker's £i£'th send will be £i.toString£ to £sends(me)(i)£.
      sends(me) = Random.shuffle(List.tabulate(iters)(i => i%n)).toArray
      for(i <- 0 until iters) dist.send(me, sends(me)(i), i.toString)
    }
    def receiver(me: Int) = thread(s"receiver $me"){
      for(i <- 0 until iters) receives(me)(i) = dist.receive(me)
    }
    run(|| (for(i <- 0 until n) yield sender(i) || receiver(i)))
    dist.shutdown()
    // Check correctness: if £r£ received £(s,m)£, then £s£ sent £m£ to £r£.
    for(r <- 0 until n; (s,m) <- receives(r)) assert(sends(s)(m.toInt) == r) 
  }
\end{scala}
\end{answerS}
